\section{Requirements Engineering}

\textbf{Mục tiêu:} Xác định các chức năng của hệ thống và những đối tượng sẽ tương tác với hệ thống Smart Event Reminder System (SERS).

\subsection{Stakeholders}
Hệ thống bao gồm các bên liên quan chính sau đây:

\begin{table}[h]
\centering
\renewcommand{\arraystretch}{1.5}
\begin{tabular}{|p{4cm}|p{10cm}|}
\hline
\textbf{Stakeholder} & \textbf{Vai trò} \\ \hline
User & Tạo và quản lý các sự kiện cá nhân. \\ \hline
Admin & Quản lý tài khoản người dùng và giám sát dữ liệu hệ thống. \\ \hline
Notification Service & Thành phần kỹ thuật chịu trách nhiệm kích hoạt và gửi nhắc nhở. \\ \hline
Project Team & Nhóm sinh viên chịu trách nhiệm thiết kế, lập trình và kiểm thử. \\ \hline
Instructor & Giảng viên đánh giá quy trình và chất lượng tài liệu kỹ thuật. \\ \hline
\end{tabular}
\caption{Danh sách các bên liên quan của dự án SERS}
\end{table}

\subsection{Yêu cầu hệ thống}

\subsubsection{Functional Requirements - FR}
\begin{enumerate}
    \item \textbf{FR1:} Hệ thống cho phép người dùng tạo sự kiện mới.
    \item \textbf{FR2:} Hệ thống cho phép người dùng nhập chi tiết sự kiện (tiêu đề, mô tả, ngày, giờ, địa điểm).
    \item \textbf{FR3:} Hệ thống cung cấp chức năng chỉnh sửa thông tin các sự kiện đã tồn tại.
    \item \textbf{FR4:} Hệ thống cho phép người dùng xóa các sự kiện khỏi danh sách quản lý.
    \item \textbf{FR5:} Hệ thống tự động gửi nhắc nhở trước khi sự kiện bắt đầu theo thời gian đã cài đặt.
    \item \textbf{FR6:} Hệ thống hiển thị thông báo nhắc nhở khi người dùng nhập sau thông tin khi tạo sự kiện.
\end{enumerate}

\subsubsection{Non-Functional Requirements - NFR}
\begin{enumerate}
    \item \textbf{NFR1 (Hiệu năng):} Thông báo nhắc nhở phải được gửi đến người dùng trong vòng tối đa 60 giây kể từ thời điểm kích hoạt.
    \item \textbf{NFR2 (Tính khả dụng):} Giao diện người dùng phải trực quan, cho phép tạo một sự kiện trong ít hơn 3 phút thao tác.
    \item \textbf{NFR3 (Độ tin cậy):} Dữ liệu sự kiện phải được lưu trữ an toàn và không bị mất khi ứng dụng khởi động lại hoặc mất kết nối mạng tạm thời.
\end{enumerate}

\subsection{User Stories và Tiêu chí chấp nhận}

\begin{longtable}{|p{1cm}|p{6cm}|p{7cm}|}
\hline
\textbf{ID} & \textbf{User Story} & \textbf{Acceptance Criteria} \\ \hline
US1 & As a User, I want to create an event with a title, date, and time so that I can schedule my plans. & 1. Hệ thống lưu sự kiện khi điền đủ thông tin bắt buộc. \newline 2. Hiển thị thông báo lỗi nếu bỏ trống tiêu đề. \\ \hline
US2 & As a User, I want to edit event details so that I can update my schedule if plans change. & 1. Thông tin mới được cập nhật ngay lập tức vào cơ sở dữ liệu. \newline 2. Người dùng có thể hủy chỉnh sửa mà không mất dữ liệu cũ. \\ \hline
US3 & As a User, I want to delete an event so that I can remove cancelled activities. & 1. Yêu cầu xác nhận trước khi xóa. \newline 2. Sự kiện biến mất khỏi danh sách sau khi xóa thành công. \\ \hline
US4 & As a User, I want to receive a reminder so that I am not late for my events. & 1. Thông báo xuất hiện đúng thời gian quy định. \newline 2. Nội dung thông báo hiển thị đúng tiêu đề sự kiện. \\ \hline
US5 & As a User, I want to receive error notifications when entering invalid data so that I can ensure the event list remains accurate. & 1. Hệ thống hiển thị thông báo lỗi cụ thể khi người dùng nhập sai định dạng dữ liệu. \newline 2. Ngăn chặn việc lưu sự kiện vào danh sách nếu có trường dữ liệu không hợp lệ. \\ \hline
\end{longtable}

\subsection{Ưu tiên yêu cầu theo mô hình MoSCoW}

\begin{table}[H]
\centering
\renewcommand{\arraystretch}{1.3}
\begin{tabular}{|l|p{10cm}|}
\hline
\textbf{Phân loại} & \textbf{Tính năng / Yêu cầu} \\ \hline
\textbf{Must have} & Tạo, sửa, xóa sự kiện; Chức năng gửi nhắc nhở tự động. \\ \hline
\textbf{Should have} & Nhập chi tiết địa điểm; Xem danh sách sự kiện theo lịch. \\ \hline
\textbf{Could have} & Mời người dùng khác; Chia sẻ sự kiện qua mạng xã hội. \\ \hline
\textbf{Won't have} & Đồng bộ hóa với Google Calendar/Outlook (trong phiên bản này). \\ \hline
\end{tabular}
\caption{Bảng phân loại mức độ ưu tiên MoSCoW}
\end{table}

\subsection{Traceability Table}

\begin{table}[H]
\centering
\renewcommand{\arraystretch}{1.3}
\begin{tabular}{|c|c|p{9.5cm}|}
\hline
\textbf{Mã FR} & \textbf{Mã User Story} & \textbf{Biểu đồ liên quan} \\ \hline
FR1, FR2 & US1 & Sequence Diagram (Create Event), Class Diagram, State Diagram\\ \hline
FR3, FR4 & US2 & Sequence Diagram (Edit/Deltete Event), Class Diagram, State Diagram \\ \hline
FR5 & US4 & State Diagram, Class Diagram \\ \hline
FR6 & US5 & Class Diagram \\ \hline
\end{tabular}
\caption{Traceability Table thể hiển liên kết giữa Yêu cầu, User Story và Diagram}
\end{table}
